\subsection{Il problema della granularità nella misurazione energetica}

Un prerequisito fondamentale per qualsiasi meccanismo di controllo
energetico è la disponibilità di misure affidabili e aggiornate a
granularità sufficientemente fine. Storicamente, la stima del consumo
energetico nei sistemi virtualizzati è stata affrontata a livello
di macchina fisica o di macchina virtuale, rendendo difficile
attribuire il costo energetico a singoli componenti applicativi.

Nel contesto FaaS, questa difficoltà è amplificata dalla natura
\emph{fine-grained} delle invocazioni e dalla variabilità delle
condizioni runtime, quali il carico concorrente, le interferenze
tra container e la distinzione tra esecuzioni a freddo
(\emph{cold start}) e a caldo (\emph{warm start}).
Negli ultimi anni, strumenti di osservabilità energetica per ambienti
containerizzati hanno reso disponibile una misurazione a granularità
significativamente più fine.

In particolare, Kepler (\emph{Kubernetes-based Efficient Power Level
Exporter}) stima il consumo energetico a livello di singolo processo
e container in ambienti Kubernetes, esportando metriche accessibili
in tempo reale tramite Prometheus \cite{Amaral2023_Kepler}.
La disponibilità di misurazioni energetiche a livello di container
costituisce un abilitatore fondamentale per integrare l'energia tra
le metriche decisionali delle piattaforme FaaS, aprendo la strada
a meccanismi di scheduling energeticamente consapevoli.
